A conversão de descrições textuais de processos em modelos BPMN 2.0 por modelos de linguagem esbarra no formato de saída: o XML da notação é extenso, e cada \textit{token} gerado é cobrado. Investiga-se se gerar uma representação intermediária compacta, expandida depois por um programa determinístico, produz modelos válidos de forma mais confiável e mais econômica que a geração direta de XML. Para isso, foram construídos uma linguagem de domínio específico, a BPMN-DSL, um transpilador que a expande em XML validado contra o esquema oficial e um modelo de sete bilhões de parâmetros ajustado com QLoRA sobre 768 pares de descrição e DSL. As abordagens foram comparadas em 46 processos de um conjunto de teste público, quanto à validade, ao custo e à fidelidade ao texto. A DSL tem cerca de 81\% menos \textit{tokens} que o XML equivalente, mas, nos mesmos modelos de grande porte, pedi-la no lugar do XML reduziu a validade de 97,8\% para 62,3\%, com a maior parte das falhas nas referências a elementos já escritos. Com a referência pelo rótulo, em análise exploratória, a validade subiu para 87,0\%. O modelo pequeno ajustado alcançou os mesmos 87,0\%, com custo por saída válida 5,5 vezes menor que o do mesmo método no modelo de grande porte e 9,4 vezes menor que o da geração direta de XML, mas com fidelidade menor, segundo uma medida exploratória. Conclui-se que a notação compacta não torna a geração mais confiável, mas a torna mais barata, e que notações destinadas a modelos de linguagem devem permitir referência ao que já foi escrito.

\textbf{Palavras-chave}: BPMN; linguagem de domínio específico; modelos de linguagem; ajuste fino; compressão de \textit{tokens}.
